\section{The Edition as a Composite of Independent Components}\label{sec:edition}

\subsection{Four adjustable components}

The separation that the LingoJam Cipher font exploits can be stated in general terms. A printed page results from a chain of choices, and four of them can be varied independently. The first is the \emph{linguistic text}, the sequence of words an author writes and a corrector repairs. The second is the \emph{character substitution}, a rule that replaces the characters of the text by other characters before anything is displayed. The third is the \emph{glyph assignment}, a rule, realised in a font, that determines which shape is drawn for each stored character. The fourth is the \emph{schedule}, a rule that determines how the second and third change across the positions of the text, and it is trivial in a static edition and central in a drifting one. Let $T\in\Sigma^{N}$ be the text, written as a sequence of $N$ letter positions, $\pi$ a substitution, and $\GG$ a space of shapes. Each stored character $x\in\Sigma$ has a \emph{route}, a path $\rho_x:[0,1]\to\GG$ from its source shape $\rho_x(0)$ to its destination shape $\rho_x(1)$, so that $\rho_x(\alpha)$ is the shape of $x$ when it has moved a fraction $\alpha$ of the way along that route. The \emph{schedule} is a vector-valued function of position,
\[
\boldsymbol{\alpha}(p)=\bigl(\alpha_{\mathrm A}(p),\alpha_{\mathrm B}(p),\ldots,\alpha_{\mathrm Z}(p)\bigr),\qquad \alpha_\ell(p)\in[0,1],
\]
where $p$ is the position in the book and $\alpha_\ell(p)$ states how far the glyph for letter $\ell$ has moved from its source form towards its destination form. The display of position $p$ is the shape $\rho_{x}\bigl(\alpha_{x}(p)\bigr)$ with $x=\pi(T_p)$, and the page is the sequence of these shapes. A \emph{uniform} transition sets every component to a common value $t(p)$, which is the single-percentage picture; a \emph{letter-by-letter} transition lets the components move separately, so that one edition may bring a single letter to its final form while leaving the others familiar, another may introduce several letters at a low fraction and strengthen them as the reader meets them repeatedly, and a third may carry every letter through ten, twenty and eventually one hundred per cent together. Figure~\ref{fig:pipeline} shows the arrangement. Two cautions attach to the notation. The value $\alpha_\ell(p)=0.5$ names a position in a specified glyph transformation, and it does not imply that the reader has acquired half the letter or that the form is halfway as difficult. And the route $\rho_x$ is itself a design choice, because the same pair of endpoints can be joined by many paths, a point developed in Section~\ref{sec:drift}.

\begin{figure}[ht]
\centering
\begin{tikzpicture}[font=\small,node distance=1.5cm,
  box/.style={draw,rounded corners=2pt,minimum height=0.9cm,minimum width=2.1cm,align=center},
  arr/.style={-{Stealth[length=2.2mm]},thick}]
\node[box] (T) {linguistic\\text $T$};
\node[box,right=of T] (E) {stored\\characters};
\node[box,right=of E] (G) {glyph\\shapes};
\node[box,right=of G] (P) {page};
\draw[arr] (T) -- node[above]{$\pi$} (E);
\draw[arr] (E) -- node[above]{$\rho_x(\alpha)$} (G);
\draw[arr] (G) -- node[above]{layout} (P);
\node[box,below=1.1cm of $(E)!0.5!(G)$,minimum width=2.6cm] (S) {schedule $\boldsymbol\alpha(p)$};
\draw[arr] (S.north) -- ++(0,0.35) -| ($(E)!0.5!(G)+(0,-0.15)$);
\draw[arr,dashed] (S.north) -- ++(0,0.35) -| ($(T)!0.5!(E)+(0,-0.15)$);
\end{tikzpicture}
\caption{The composite of an edition. The text, the substitution $\pi$, the routes $\rho_x$ and the schedule $\boldsymbol\alpha$ are separately identifiable, and the page depends on all four.}\label{fig:pipeline}
\end{figure}

Two configurations at the extremes show what the separation permits. In a \emph{compensated} configuration the stored text is substituted and the glyph assignment $\gamma:\Sigma\to\GG$, a static map from stored characters to shapes (a route evaluated at a fixed fraction), undoes the substitution, so that $\gamma(\pi(x))$ is the ordinary shape of $x$. The page looks entirely normal, and the stored characters are the cipher. This is the arrangement of the LingoJam Cipher font, where the stored characters belong to one alphabet and the font displays them as another. In an \emph{uncompensated drifting} configuration the stored text is ordinary and $\pi$ is the identity, while the routes carry each letter from its ordinary shape to the shape of its partner as the schedule advances. Here the stored characters are the plain text, and the visible forms change. Between them lie intermediate arrangements in which the font retains recognisable features of both letters of a pair, so that the substitution remains perceptible while the familiar structure supports reading, and one can also produce several grades of the same idea, moving from strongly corrective glyphs towards ordinary encoded letters, so that the reader's familiarity with the alphabet gradually takes over the work that the font initially performed.

\subsection{Display equivalence and the storage decision}

A given page can be produced by many different pairs of stored text and glyph assignment. If $X$ is a stored string and $\gamma$ a glyph assignment, the display is $\gamma(X)$, and two pairs $(X,\gamma)$ and $(X',\gamma')$ produce the same page precisely when $\gamma(X)=\gamma'(X')$. For instance, with the exchange of B and P, the pair (stored PUB, assignment that draws P as B and B as P) and the pair (stored BUP, ordinary assignment) look identical on the page. Which representative of the equivalence class is stored is a decision with consequences, because the stored text is what search, copying, screen readers, speech synthesis, spell checking and version control act on. A compensated cipher stores gibberish that merely looks like text, and every such tool sees the gibberish. An uncompensated drifting configuration stores the author's text, so that these tools continue to work, and places the whole burden of the transformation on the glyphs. This suggests a design principle, which will be called the stored-text principle: wherever possible the stored text should be the linguistic text, and transformations should be carried by the glyph assignment and the schedule. The principle is not a theorem but a way of making accessibility and correction survive the experiment. It means in particular that a plain edition can always be produced from the same source by selecting the ordinary glyph assignment, and that a reader who finds the drifting edition unreadable is never denied access to the text.

\subsection{Font technology}

The relevant mechanisms are standard but their limits matter. In an OpenType font the character map associates each stored character with a glyph, and substitution tables in the glyph substitution table allow replacement of glyphs by alternates, in some cases depending on neighbouring glyphs \citep{opentype}. A font whose only job is to display each character as another character's shape is a simple case. Interpolation between shapes is available through variable fonts, introduced in OpenType 1.8, in which a font defines design axes and glyph outlines vary continuously with the axis values, provided that the master outlines are point-compatible, meaning that they have the same structure of contours and points so that corresponding points can be interpolated \citep{opentype}. The idea of a parametrised letter is older than the standard; Knuth's description of a meta-font treats a typeface as a program with parameters from which many related designs can be generated \citep{knuth1982}.

Two limitations determine how a drifting book is built. First, a font's layout rules operate on the text of a run and do not count occurrences across a book, so the stage of a letter, which depends on how many times that letter has been seen, cannot be computed inside the font. It has to be assigned by a preprocessing step that annotates each token with its stage, either by selecting among several static fonts that each embody one stage, or by setting the axis values of a variable font for each run of text. Second, an interpolation axis in a variable font is applied to every glyph in a run in the same way, so that per-letter schedules require either separate axes per letter, which becomes cumbersome for an alphabet, or a division of the text into runs, one for each letter and stage, at the cost of heavy markup. The sensible organisation is an edition builder, a program that reads the text, computes stages from the schedule, emits annotated text, and calls a typesetting engine with the appropriate fonts. The same builder can produce the plain edition, a static cipher, and a drifting edition from one source, which is what the reproducibility requirements of Section~\ref{sec:source} demand.

\subsection{Contextual display}

The letterwise model, in which each stored character receives a shape independent of its neighbours, is adequate for the Latin alphabet in its ordinary use and fails for scripts whose letter forms depend on context. Arabic is the standard example, and Section~\ref{sec:arabic} treats it. The general model is that the display function is not a map $\Sigma\to\GG$ extended letter by letter but a map from strings of stored characters to strings of glyphs, in the simplest case computed by a finite-state transducer reading the string and emitting glyphs, possibly with ligatures that merge several characters into one glyph. All the statements of this section go through with this generalisation, but the commuting of rendering with substitution, which is automatic in the letterwise case, is not, and Proposition~\ref{prop:commute} below records exactly when it holds.
